\chapter{CONCEITOS FUNDAMENTAIS} \label{conceitos}

Este capítulo apresenta os conceitos usados ao longo do trabalho, a partir das palavras-chave do resumo e dos termos empregados nos trabalhos relacionados. Para cada conceito, são apresentados sua definição, sua origem, sua finalidade e suas limitações.

\section{Grandes Modelos de Linguagem}

Grandes Modelos de Linguagem (LLMs) são modelos treinados em grandes volumes de texto para gerar respostas em linguagem natural. A base da maioria desses modelos é a arquitetura Transformer, proposta por \textcite{vaswani2017attention}, que usa mecanismos de atenção para relacionar as palavras de um texto. \textcite{brown2020language} mostraram, com o GPT-3, que ampliar a escala desses modelos permite que eles resolvam tarefas variadas a partir de poucas instruções ou exemplos.

Hoje os LLMs já mediam decisões e a geração de conteúdo em diversas atividades. Sua principal limitação, para este trabalho, é que eles geram texto fluente sem verificar fatos: uma resposta bem escrita não é, necessariamente, uma resposta correta, e o modelo pode apresentar informações falsas com a mesma segurança de informações verdadeiras.

\section{Alinhamento por preferências}

O alinhamento por preferências é a etapa em que o modelo é ajustado para produzir respostas bem avaliadas por pessoas ou por modelos de preferência. A técnica mais conhecida é o RLHF (\textit{Reinforcement Learning from Human Feedback}), proposta por \textcite{christiano2017deep} e aplicada a modelos de linguagem por \textcite{ouyang2022training}, que a usaram para fazer o modelo seguir instruções de forma útil e cooperativa.

A limitação dessa técnica é que ela otimiza o que os avaliadores preferem, e não necessariamente o que é verdadeiro. Como as avaliações às vezes favorecem respostas que concordam com o usuário em vez de respostas corretas, o processo pode criar uma tensão entre agradar o usuário e preservar a veracidade \cite{sharma2024towards}.

\section{Sycophancy}

\textit{Sycophancy} é a tendência de um modelo acomodar a crença ou a sugestão do usuário mesmo quando ela é incorreta. \textcite{perez2023discovering} fizeram uma das primeiras medições sistemáticas do fenômeno em LLMs, e \textcite{sharma2024towards} o relacionaram ao alinhamento por preferências. Em diálogos com múltiplos turnos, a pressão sustentada do usuário pode levar o modelo a abandonar uma resposta correta e mudar de posição \cite{hong2025measuring}.

Medir a \textit{sycophancy} serve para avaliar a confiabilidade de um modelo quando o usuário discorda dele. A principal limitação do conceito é que nem toda mudança de posição é sicofântica: se o usuário aponta um erro real, mudar de resposta é o comportamento correto. Por isso, o fenômeno precisa ser medido junto com a correção seletiva.

\section{Capitulação e resistência}

Diante da pressão do usuário, o modelo pode ter dois comportamentos. Ele \textbf{resiste} quando mantém a evidência e corrige a premissa incorreta do usuário. Ele \textbf{capitula} quando muda de posição e valida a sugestão sem base. \textcite{hong2025measuring} caracterizaram a capitulação em diálogos multi-turno pelo turno em que ocorre a primeira mudança de posição (\textit{Turn of Flip}) e pelo número de mudanças ao longo do diálogo (\textit{Number of Flip}).

Essas medidas mostram não só se o modelo cede, mas quanto tempo ele resiste. Sua limitação é depender do critério usado para decidir o que conta como mudança de posição, especialmente quando o modelo responde de forma ambígua.

\section{Correção seletiva}

A correção seletiva é a capacidade de o modelo aceitar sugestões válidas e rejeitar sugestões inválidas, separando a correção legítima da cessão indevida à pressão. \textcite{sinha2026sycobench} propôs medir esses dois comportamentos separadamente, já que um modelo que nunca muda de posição não é sicofântico, mas também não aceita correções.

Para ser medida, a correção seletiva exige saber com certeza qual é a resposta correta de cada questão, o que limita sua aplicação a perguntas com resposta objetiva e verificável.

\section{Criticidade e gravidade potencial}

A criticidade indica o potencial de dano de uma resposta incorreta em um determinado domínio. A mesma concordância indevida pode ser um erro inofensivo em um domínio e uma orientação potencialmente prejudicial em outro. Neste trabalho, propõe-se uma definição própria de criticidade, inspirada na noção de gravidade usada em matrizes de risco. A classificação abaixo é preliminar e depende da avaliação do grupo e do orientador, conforme discutido na metodologia:

\begin{itemize}
    \item \textbf{Baixa}: o erro é inofensivo (por exemplo, conhecimentos gerais);
    \item \textbf{Média}: o erro pode causar prejuízo financeiro ou de direitos (por exemplo, finanças pessoais);
    \item \textbf{Alta}: o erro pode causar risco físico (por exemplo, saúde e segurança).
\end{itemize}

A limitação dessa classificação é ser subjetiva. Para reduzir esse efeito, propõe-se que dois integrantes revisem, de forma independente, a criticidade de cada cenário e discutam as divergências. Esse procedimento ainda depende de aprovação; a classificação não deve ser determinada apenas pelo domínio.
